\documentclass[answers,12pt,addpoints]{exam} 
\usepackage{import}

\import{C:/Users/prana/OneDrive/Desktop/MathNotes}{style.tex}

% Header
\newcommand{\name}{Pranav Tikkawar}
\newcommand{\course}{Spatial Statistics}
\newcommand{\assignment}{Homework 3.2}
\author{\name}
\title{\course \ - \assignment}

\begin{document}
\maketitle

\begin{questions}
\question The concept of a microstructure in spatial statistics is a really cool idea. I usually don't about how measurement error can affect the data and when I do it is usually in the content of IID normal measurement error. I really want to learn more about how non-normal measurement error can affect the data and how to model it. 

\question I find the averaging of $Z$ given some integrable function $v$ to be a great idea for smoothing out the data as well as for making measurments in a more continuous way. I am curious how choice of $v$ can affect the data and how to choose $v$ in practice since there needs to be a balance between smoothing the data and keeping the data close to the original.

\question I am having difficulty understanding the concept of the bessel functions and how they are used in the context of spatial statistics. I am aware they are derived from the differential equation $x^2y'' + xy' + (x^2 - \nu^2)y = 0$ and there exists many different types of bessel functions. What are the modified bessel functions and where are they derived from? 

\question I am failing to see intuitively how the PSD property and the monotone property are related. I think that it stems from the interpretaion of the $\int_0^\infty e^{-tu}dG(u)$ function. This looks like a Laplace transform and I think this is the connection between the two properties but not sure how to make the connection more clear.

\question It is unfortunate that the Cauchy model is not fit for application for spatial data. The model is simple but the differentiability issues make it clear that it is not a good model for spatial data. I was curious if there were any datasets that might come to mind where the Cauchy model might be a good fit.

\question I like this combination of models using the nice long range properties of the Cauchy model and the nice short range properties of the Matern model to create a model that is good for both short and long range dependence. I do not understand why 6 is a large amount of parameters for a model. Is there an issue when it comes to estimation or is it just that it is more difficult to interpret the model with more parameters?

\question I am frankly very lost for generalized autocovariance functions and thier uses. This chapter 4.6 is very information dense and I intend on reading more tonight to make sure I understand the material. 

\question Overall, this days reading was dense and I am not sure I understand all of the material. I will be rereading the material and potentially will bring questions for class. I find the concepts interesting and relevant, but I am having trouble understanding the material in depth. 

\end{questions}

\end{document}